\cotefile{157-5}
\scanpage{52}
\mspage{39}
% archive sheet number: 51

3. \uline{Catégories de modèles à \unsure{cofibrations|cofibration}}

Soit $M$ une catégorie munie de deux \unsure{ens} de flèches
\[
 W,C\subset\mathrm{Fl}\,M,\unsure{.}
\]

On suppose que $M$ contient un objet initial, \unsure{et}

\deleted{\illegible{1}} $\mathrm{Cof}_{0}$ : $W$ est \unsure{modérément} saturé i.e. contient iso et si $f,g$ composables et si deux parmi $f,g,gf$ sont dans $W$, le troisième aussi

$\mathrm{Cof}_{1}$ : $C$ \deleted{\unsure{stable}}contient iso, stable par composition et par \unsure{cochangement} \unsure{des bases|de base}

$\mathrm{Cof}_{2}$ : $TC\overset{\mathrm{def}}{=}C\cap W$ est stable par cochangement de base

$\mathrm{Cof}_{3}$ : Toute \unsure{flèche|flèches} $f$ dans $M$ s'écrit
\[
 f=pi
\]
avec $i\in C,\ p\in W$.

\inserted{Ce sont les} Axiomes de K.~S. Brown \inserted{(\unsure{forme duale})}, qui suppose de plus que les objets de $M$ sont $C$-cofibrants, i.e. $\varnothing_M\unsure{\mathbin{-}|\longrightarrow} X$ est dans $C$. Je ne ferai pas cette hypothèse. La sous-catégorie $M_C$ \deleted{\unsure{de}} de $M$ formée des
\scanpage{53}
\mspage{40}
% archive sheet number: 52

objets $C$-cofibrants satisfait aux axiomes de Brown, de plus on \deleted{\unsure{voit}} \inserted{\unsure{prouve}} que
\[
 \mathrm{Ho}_{W}(M_C)\xrightarrow{\approx}\mathrm{Ho}_{W}(M)
\]

est une équivalence. Ainsi les résultats de Brown relatifs à $M_C$, \unsure{muni|munie} des \unsure{ens} de flèches induits par $W,C$, sont applicables à $\mathrm{Ho}_W(M)$ et à ses avatars (tels $\mathrm{Ho}_W({}_S\backslash M)$ ; $\mathrm{Ho}_{\unsure{W|\infty}}(\deleted{\unsure{{}_S\backslash}}M/_{T})$ et $\mathrm{Ho}_{W}({}_S\backslash M_{u_0}/_{T})$ où $u_0:S\unsure{\mathbin{-}|\longrightarrow} T$ est donné.

Il se trouve que pour l'essentiel, les résultats homotopiques \unsure{standart|standarts|standards} (notamment suite \unsure{coexte|coexacte} de cofibration) relatifs aux opérations \unsure{homotopiques} ``à droite'' (ou faut-il dire à gauche ? il s'agit \unsure{pourtant|précisément} des questions d'exactitude à droite, non à gauche) sont valables. Il \deleted{\unsure{me}} faudra revenir dessus de façon plus circonstanciée par la suite. \deleted{Je suis \unsure{particulièrement}} \inserted{Notamment il y} a une bonne théorie des carrés $W$-cocartésiens (et \uline{dans $M_C$}, il y a assez de $W$-cofibrations$^{\times}$, \ldots).

% margin note written sideways
\margin{$^{\times}$plus précisément, par K.S. Brown, \uline{dans} $M_C$ les objets de $C$ sont des $W_C$-cofibrations.}
\scanpage{54}
\mspage{41}
% archive sheet number: 53

et de plus, ceux-ci sont \uline{\unsure{exacts}}. La catégorie $\mathrm{Ho}_W(M)\approx\mathrm{Ho}_W(M_C)$ ``se calcule'' par un calcul de fractions à gauche dans $\pi M_C$, où $\pi M_C$ est la catégorie déduite de $M_C$ en y remplaçant les $\mathrm{Hom}_{M_C}(X,Y)=\mathrm{Hom}_M(X,Y)$ par le quotient par une relation d'homotopie convenable :

$f,g:X\unsure{\rightrightarrows|=}Y$ sont \unsure{QB-homotopes} (au sens de \uline{Quillen}-\uline{Brown}) s'il existe un $X'\xrightarrow{\alpha}X$ dans $W$ (\deleted{\unsure{et $X'$}} $X'\in\mathrm{Ob}\,M_C$ également) tel que $f\alpha\overset{\unsure{l}}{\sim}g\alpha$ au sens de Quillen, i.e. existe un objet-cylindre

\diagram{\begin{tikzcd}[ampersand replacement=\&,row sep=small,column sep=small]
 I_{X'} \arrow[r,"{\in W}"] \& X' \\
 X'\vee X' \arrow[u,"{\in C}"] \arrow[ur,dashed] \&
\end{tikzcd}}

pour $X'$, tel que $(f\alpha,g\alpha):X'\vee X'\longrightarrow Y$ se factorise par un $h':I_{\unsure{X''|X'}}\longrightarrow Y$.

C'est de ce critère que résulte l'exactitude \inserted{(\unsure{à g.|à Q.|à D.})} des carrés $W$-cocartésiens.

Je veux voir comment une structure de cofibrations s'étend aux catégories $M(I^{\unsure{\circ|0}})=\underline{\mathrm{Hom}}(I,M)$. \deleted{On suppose d'abord} \inserted{On a dans $M(I^\circ)$}
\scanpage{55}
\mspage{42}
% archive sheet number: 54

\deleted{$I$ finie} l'ens $W_{I^{\unsure{\circ|0}}}$, et $C_I$ (il vaut mieux sans doute le noter aussi $C_{I^\circ}$) \unsure{défini|définis} au \S{} précédent. Il s'agit de vérifier les axiomes. L'axiome $\mathrm{Cof}_0$ est immédiat, et bien connu pour $I$ une catégorie quelconque.

$\mathrm{Cof}_1$ : $C_{I^\circ}$ contient les iso, c'est clair, parce que $C$ les contient. Stable par composition $X_{\unsure{\cdot|\bullet}}\xrightarrow{u}Y_{\cdot}\xrightarrow{v}Z_{\cdot}$, \unsure{regardons} $K\subset J$, posons $w=\unsure{v\circ u|vu}$

\diagram{\begin{tikzcd}[ampersand replacement=\&,column sep=small,row sep=small]
 X_K \arrow[r] \arrow[d,dash] \& X_{\deleted{\unsure{K}}J} \arrow[d] \& \& \\
 Y_K \arrow[r] \arrow[d] \& Y_{K,J} \arrow[r,"{u_{K,J}}"] \arrow[d] \& Y_J \arrow[d,dash] \& \\
 Z_K \arrow[r] \& \cdot \arrow[r,"{u'_{K,J}}"] \arrow[rr,phantom,shift right=2ex,"{\underbrace{\hspace{8em}}_{w_{K,J}}}"'] \& Z_{K,J} \arrow[r,"{v_{K,J}}"] \& Z_J
\end{tikzcd}}

$w_{K,J}$ est composé de $v_{K,J}$, qui est dans $C$, et $u'_{K,J}$, qui l'est car se déduit de $u_{K,J}$ par \unsure{cochgt} de base. (NB \unsure{On} a déjà fait le raisonnement avec un seul $u$, et \uline{trois} ouverts $J\subset K\subset L$. Il
\scanpage{56}
\mspage{43}
% archive sheet number: 55

y aurait lieu d'introduire la notion de $C$-carré

\diagram{\begin{tikzcd}[ampersand replacement=\&,column sep=small,row sep=tiny]
 X_0\arrow[r,dash]\arrow[d,dash] \& X_1\arrow[d,dash]\\
 X_{\unsure{2|\dot{2}}}\arrow[r,dash] \& \unsure{x|X}_3
\end{tikzcd}}

, comme un carré tel que $X_{\unsure{1}}\amalg_{X_0}X_2\unsure{\mathbin{-}|\longrightarrow}X_3$ soit dans $C$\inserted{(NB notion symétrique pour les \unsure{opérations|foncteurs} de symétrie des carrés)}, et dire que le composé de deux $C$-carrés est un $C$-carré. $]$

Cochangement de base

\diagram{\begin{tikzcd}[ampersand replacement=\&,row sep=small]
 X_{\unsure{\cdot|\bullet}}\arrow[r,"u"]\arrow[d] \& Y_{\cdot}\arrow[d,dash]\\
 X'_{\cdot}\arrow[r,dash,"u'"] \& Y'_{\cdot}
\end{tikzcd}}

, Considérons $\deleted{\unsure{J}}K\subset J$, on a

\diagram{\begin{tikzcd}[ampersand replacement=\&,column sep=small,row sep=small]
 \& X_K\arrow[r]\arrow[d,dash]\arrow[ddl,thick] \& X_{\deleted{\unsure{K}}J}\arrow[d]\arrow[dr,thick]\arrow[ddl,thick] \&\\
 \& Y_K\arrow[r]\arrow[ddl]\arrow[ur,phantom,"{\text{\unsure{cocart}}}" description] \& Y_{K,J}\arrow[r,dash,"{u_{K,J}}"']\arrow[ddl] \& Y_J\arrow[ddl,thick]\\
 X'_K\arrow[r,leftrightarrow,thick]\arrow[d,thick]\arrow[ur,phantom,"{\text{cocart}}" description,sloped] \& X'_J\arrow[d,thick]\arrow[dr,thick] \& \&\\
 Y'_K\arrow[r,leftrightarrow,thick]\arrow[ur,phantom,"{\text{cocart}}" description] \& Y'_{J,K}\arrow[r,thick,"{u'_{\unsure{K},J}}"] \& Y'_J\arrow[uur,phantom,"{\text{cocart}}" description,sloped] \&
\end{tikzcd}}

, où les trois carrés \inserted{\unsure{latéraux}} marqués ``cocart'' sont cocartésiens, ainsi que le carré en plan incliné. Il en résulte, d'abord que le quatrième carré latéral du cube est lui aussi cocartésien,
\scanpage{57}
\mspage{\unsure{44|41}}
% archive sheet number: 56

puis que le carré ayant deux côtés $u_{K,J},u'_{K,J}$ est cocartésien. (NB Comparer le diagramme similaire p.~32 -- il y aurait un lemme commun à dégager, distinguant \unsure{deux|des} cas \ldots)

\deleted{Pour prouver que} ($\mathrm{Cof}_{\unsure{2|1}}$) Pour prouver que \deleted{$\unsure{f}$} $C_{I^{\unsure{\circ|0}}}\cap W_{I^\circ}$ est stable par cochangement de base, je ne vois pas de façon de le faire, si je ne sais déjà que

\[
 C_{I^\circ}\cap W_{I^\circ}=(TC)_{I^\circ},\tag*{(\unsure{\alpha})}
\]

et alors on applique \unsure{le} \deleted{\unsure{vérificat}} résultat précédent (où $W$ n'entrait pas) à $TC$ \unsure{ald} $C$. Mais pour avoir \unsure{ceci|ces}, il faut que je suppose qu'on est dans l'un des deux cas suivants (de façon à pouvoir appliquer le th.~2.1. et son corollaire 2.5.)

a) $I$ fini, ou plus généralement $I$ \unsure{artinien} \uline{et} $\mathrm{Top}(I)$ est noethérien

b) $I$ artinien, et de plus $C$ et $TC$ stables par composition transfinie.

($\mathrm{Cof}_3$) L'axiome de factorisation me paraît \unsure{à|a} priori le plus délicat$^{\times}$. Mais cela \deleted{\unsure{tout}} fonctionne sans trop de mal. Soit donné
\[
 f:X_{\unsure{\cdot|\bullet}}\unsure{\mathbin{-}|\longrightarrow}Y_{\cdot}.
\]

% slanted
\margin{$^{\times}$L'argument qui suit s'applique \unsure{à d'autres} \unsure{situations} de factorisation, par exemple pour une str. \unsure{proprement} de Quillen, donnée $(\unsure{\Phi,\Psi|C,F})$ dans $M$.}
\scanpage{58}
\mspage{45}
% archive sheet number: 57

dans $M(I^{\unsure{\circ|0}})$, considérons les \unsure{triples|4-tuples|k-tuples}
\[
 (J,Z_{\deleted{\unsure{J}}},\deleted{\unsure{p}}i_{\deleted{\unsure{J}}},p_{\deleted{\unsure{J}}})
\]

où $J\subset I$ est un ouvert, $Z_{\deleted{\unsure{J}}}$ un objet de $M(J^\circ)$, et $i_{\deleted{\unsure{J}}},p_{\deleted{\unsure{J}}}$ sont \unsure{des|les} flèches
\[
 \left\{\begin{aligned}
 &X|J\xrightarrow{i} Z\xrightarrow{p}Y|{\unsure{J}}\\
 &\text{avec}\quad i\in C_{J^\circ},\quad p\in W_{J^\circ}.
 \end{aligned}\right.
\]

On ordonne les quadruples par prolongement. Si on a un ens. \unsure{tot.} ordonné, de $(J_\alpha,Z_\alpha,\deleted{\unsure{f}}i_\alpha,p_\alpha)$, posant $J=\bigcup J_\alpha$, on trouve des $Z,i,p$ \deleted{\unsure{qui}} \inserted{uniques} \unsure{sur} $J$ qui prolongent les $Z_\alpha,i_\alpha,p_\alpha$ donnés, de plus on aura $i\in C_{J^\circ},p\in W_{J^\circ}$, car l'appartenance à \deleted{\unsure{des classes}} $C_{J^\circ}$ resp. $W_{J^\circ}$ est question locale sur $J$. Donc \inserted{par Zorn,} on est ramené à prouver qu'\deleted{il y a un} si $(J,Z,i,p)$ est maximal, on a $J=I$ -- d'où la factorisation voulue $f=pi$.

Donc : prouver que si $J\ne I$, le quadruple n'est pas maximal. Considérons le fermé \inserted{non vide} $I\setminus J$ \deleted{$\unsure{=A}$}. $I$ est supposé artinien, il y a un \uline{élément minimal} $\unsure{i_0}$ de $I\setminus J$. Posons
\scanpage{59}
\mspage{46}
% archive sheet number: 58

\[
 J'=J\cup\{i_0\}=J\cup I_{/i_0}
\]

et montrons qu'on peut prolonger le \deleted{\unsure{quadru}} \inserted{triplet} $Z,i,p$ en $Z',i',p'$ sur $J'$. Posant $\deleted{\unsure{I}}K'=I_{/i_0}$, $K=J\cap K$\inserted{$=I_{<i_0}$}, on a donc

\diagram{\begin{tikzcd}[ampersand replacement=\&,column sep=large]
 K=J\cap K'\arrow[r,hook]\arrow[d,hook] \& K'=I_{/i_0}\arrow[d,hook]\\
 J\arrow[r,hook] \& J'=J\cup K'
\end{tikzcd}}

On voit donc que les prolongements de $(Z,i,p)$ de $J$ à $J'$ correspondent à ceux de $(Z|K',i|K',p|K')$ à $K$. Pour ceci, on note que pour une flèche $g$ \deleted{\unsure{d'objets}} objets de $M(J'^{\unsure{\circ|0}})$, cette flèche est dans $C_{J'^\circ}$ resp. $W_{J'^\circ}$, ssi ses restrictions à $J$ et à $K'$ le sont.

On est donc ramené au cas \unsure{standart|standard} d'un ensemble ordonné \deleted{\unsure{$K$ ayant}} avec objet final $\unsure{e_0|i_0}$, ens. ordonné qu'on va noter $I$, et l'ens ordonné $I'$ des objets $i<e_0$, i.e. $I'=I\setminus\{e_0\}$, on a $f:X_{\unsure{\cdot|\bullet}}\to Y_{\cdot}$ \deleted{\unsure{flèche}} dans $M(I^\circ)$, une factorisation de $f|I'$ en $p'\circ i'$ avec $i'\in C_{I'^\circ}$, $p'\in W_{I'^\circ}$,
\scanpage{60}
\mspage{47}
% archive sheet number: 59

et il s'agit de la prolonger en une factorisation similaire de $f$. On note que la donnée d'un objet $X_{\unsure{\cdot|\bullet}}$ de $M(I^{\unsure{\circ|0}})$ équivaut à celle d'un objet $X'_{\cdot}$ de $M(I'^\circ)$, (la restriction \inserted{de $X_{\cdot}$} à $I'$) et d'une flèche
\[
 X'_{I'}\xrightarrow{\alpha_X}X_{\unsure{e_0|i_0}}=X_0
\]

dans $M$, où $X'_{I'}=\underset{i'\in I'}{\underline{\lim}}\,X'_{i'}$. Donc la donnée de la flèche $f:X_{\cdot}\unsure{\mathbin{-}|\to}Y_{\cdot}$ revient à celle de $f'=f|I':X'_{\cdot}\mathbin{-}Y'_{\cdot}$, et d'une flèche $f_{\deleted{\unsure{i_0}}0}:X_{\deleted{\unsure{i_0}}0}\mathbin{-}Y_{i_0}$ rendant commutatif le carré

\diagram{\begin{tikzcd}[ampersand replacement=\&]
 X'_{I'}\arrow[r,"\alpha_X"]\arrow[d,"{f'_{I'}}"'] \& X_{\deleted{\unsure{i}}0}\arrow[d,"{f_{\unsure{0|i_0}}}"]\\
 Y'_{I'}\arrow[r,dash,"\alpha_Y"] \& Y_{\unsure{0|i_0}}
\end{tikzcd}}

De $=$, trouver une factorisation de $f$ en $pi$, qui prolonge \deleted{\unsure{la}}une factorisation donnée $f'=p'i'$ sur $J'$, revient à la donnée d'un diagramme commutatif

\diagram{\begin{tikzcd}[ampersand replacement=\&,column sep=small,row sep=small]
 \& \& X_{I'}\arrow[rr,"\alpha_X"]\arrow[dl,"{i'_{I'}}"']\arrow[dd,"{f'_{I'}}"] \& \& X_{\unsure{0_0|i_0}}\arrow[dl,"{i_{\unsure{0|i_0}}}"']\arrow[dd,"{f_{\unsure{i_0|0}}}"]\\
 (\unsure{M|\alpha|\ast} \& Z_{I'}\arrow[rr,"\alpha_Z"]\arrow[dr,"{p'_{I'}}"'] \& \& Z_{\unsure{0_0|i_0}}\arrow[dr,"{p_0}"'] \&\\
 \& \& Y_{I'}\arrow[rr,"\alpha_Y"] \& \& Y_{\unsure{0|i_0}}\ ,
\end{tikzcd}}
\scanpage{62}
\mspage{\unsure{48}}
% archive sheet number: 61

où l'objet $Z_{\unsure{i_0|0}}$, et les trois flèches $\alpha_Z,i_0,p_0$ sont à déterminer. Si on a $i'\in C_{I'^{\unsure{\circ|0}}},p\in W_{I'^\circ}$, alors on conclut par th.~2.1 $1^\circ$ (\unsure{cor.} 2.2) que \inserted{$i'_{I'}\in C$,} et si on veut que $i\in C_{I^\circ},p\in W_{I^\circ}$, il reste seulement à exiger \deleted{d'une part que}

\[
 p_0\in W,\ (i'_{I',I}:\underbrace{Z_{I'}\amalg_{X_{I'}}X_0}_{Z'_0}\longrightarrow Z_0)\in C
\]

d'après le th.~2.1 $2^\circ$. On va alors construire le diagramme $(*)$, à partir du diagramme similaire ci-contre (avec $Z_0$ de $(*)$ remplacé par $Z'_0=Z_{I',I}$)

\diagram{\begin{tikzcd}[ampersand replacement=\&,column sep=small,row sep=small]
 \& X_{I'}\arrow[rr,"\alpha_X"]\arrow[dl,"{i'_{I'}\in C}"']\arrow[dd,"{f'_{I'}}"] \& \& X_0\arrow[dl,"{i'_0}"']\arrow[dd,"{f_0}"]\arrow[dr,dashed,"{i_0=j_0i'_0}"] \&\\
 Z_{I'}\arrow[rr,"{\alpha'_Z}"]\arrow[dr,"{p'_{I'}}"']\arrow[urrr,dotted,leftrightarrow,shorten <=1.6em,shorten >=2.7em,"{\text{cocart}}" description,sloped] \& \& Z'_0=Z_{I',I}\arrow[rr,dash,dotted,very thick]\arrow[rr,thick,"{j_0\in C}"]\arrow[dr,"{p'_0}"'] \& \& Z_0\arrow[dl,dashed,"{p_0\in W}"]\\
 \& Y_{I'}\arrow[rr,"\alpha_Y"] \& \& Y_0 \&
\end{tikzcd}}

en factorisant la flèche canonique $p'_0:Z'_0\longrightarrow Y_0$ en \deleted{$\unsure{p}$}
\[
 p'_0=\deleted{\unsure{q}}p_0j_0,\quad\text{avec}\quad \deleted{\unsure{q}}p_0\in W,\ j_0\in C,
\]

et en posant
\[
 i_0=j_0i'_0,
\]
ce qui achève la construction.

On a ainsi prouvé le
\scanpage{63}
\mspage{50}
% archive sheet number: 62

\uline{Théorème} 3.1. Soit $(M,W,C)$ une catégorie à cofibrations (cf page 39), \inserted{$M$ stable par $\unsure{\varinjlim|\lim}$ finies.} Pour tout ensemble ordonné $I$, considérons le triple $(M(I^{\unsure{\circ|0}}),W_{I^\circ},C_{I^\circ})$ correspondant, où $C_{I^\circ}$ est défini au \S{}2 (cf à ce sujet le th.~2.1). Supposons que \deleted{\unsure{soit}} $I$ soit un ens ordonné \uline{artinien}, (p.\ ex. $I$ \unsure{loc.}$^{\unsure{t|mt}}$ fini i.e. les $I_{/i}$ finis). Supposons de plus qu'on soit dans l'un des deux cas suivants

a) L'espace $\unsure{\mathrm{Top}}(I)$ est noethérien i.e. toute suite croissante d'ouverts de $I$ est stationnaire (p.\ ex. $I$ fini).

b) $C$ et $TC=C\cap W$ sont stables, non seulement par composition \unsure{ordinaires|ordinaire} mais aussi par compositions transfinies. Dans ce cas, on suppose aussi $M$ stable par \inserted{petites} $\varinjlim$ \deleted{``quelconques''} filtrantes (donc par \inserted{petites} $\varinjlim$ quelconques).

Sous ces conditions, $(M(I^\circ),W_{I^\circ},C_{I^\circ})$ est également une catégorie à cofibrations.

\uline{\deleted{\unsure{Corollaire}}} \inserted{Prop\unsure{.}} 3.2. Soit $u:I\longrightarrow I'$ une application croissante entre ensembles ordonnés, d'où un couple de foncteurs adjoints

\diagram{\begin{tikzcd}[ampersand replacement=\&,column sep=large]
 (3.\unsure{2|1}.1) \& M(I^\circ)\arrow[r,shift left=0.6ex,"{u_!}"] \& M(I'^\circ)\arrow[l,shift left=0.6ex,"{u^*}"]
\end{tikzcd}}
\scanpage{64}
\mspage{51}
% archive sheet number: 63

a) Le foncteur $u^{\unsure{*|\alpha}}$ est compatible avec les localiseurs $W_{I'^{\unsure{\circ|0}}},W_{I^\circ}$. Si $I$ est \unsure{discrèt|discrète|discret}, ou $u$ étale, alors $u^*$ applique $C_{I'^\circ}$ dans $C_{I^\circ}$.

\deleted{[Il en est encore de $=$ si $u$ est \deleted{\unsure{co}}fibrant, \deleted{\unsure{ses}} si \deleted{chaque} \inserted{\unsure{de plus}} \deleted{fibre est un objet initial}, pour \deleted{tt} $s'\in I'$, \deleted{les} \inserted{toute} composante connexe de $I_{s'}$ \deleted{est} a un objet initial, et si de plus, pour \deleted{tt flèche} \inserted{\unsure{flèches}} $s'\le t'$ dans $I'$, le foncteur d'image inverse $I_{t'}\longrightarrow I_{s'}$ transforme objet initial d'une composante connexe en objet initial de la composante connexe correspondante.]}

Il en est \unsure{même} de $=$ quand $u$ est loc. \deleted{\unsure{fort.}} cofiltrant, i.e si pour tt \deleted{$\unsure{g}$} $i'\in\mathrm{Ob}\,I'$, la catégorie $i'\backslash I$ est loc. \deleted{\unsure{fort.}} \unsure{filtr.} \deleted{\unsure{et}}, i.e. \deleted{les catégories localisées $(i'\backslash I)$, $u$ est} \inserted{ses composantes connexes sont} \deleted{\unsure{pouvant}} \deleted{\unsure{filt}} cofiltrantes. (condition $\mathrm{PS}_g\,1$ sur $i'\backslash I$).

% slanted
\margin{inutile de parler de ``\unsure{fort}'' loc. cofiltrant, car il s'agit \unsure{d'ens.|d'un} \unsure{(pré)ordonnés|(co)préordonné}, où il n'y a pas \unsure{de distinction|à distinguer} entre notions \unsure{forte} et \unsure{faible}.}

b) \deleted{Supposons $I$ artinien, et qu'on soit dans l'un des deux cas a), b) du th.~3.1.} \inserted{Dans tous les cas,} \deleted{Alors} $u_!$ transforme \deleted{$M(I^\circ)_C$ en \unsure{$M($}} $C_{I^\circ}$ en $C_{I'^\circ}$, $TC_{I^\circ}$ en $TC'_{I'^\circ}$, enfin $M(I^\circ)_C$ en $M(I'^\circ)_C$.

% slanted
\margin{\deleted{\unsure{Ici, inutile} que $I'$ soit un \unsure{ens.} ordonné}}

\uline{Démonstration} Prouvons a), dans le cas le plus général, où on suppose seulement que $u$ est loc. \unsure{fort|ft} cofiltrant. On applique II 7.\unsure{4} (p.~\unsure{92|97}), qui nous dit que $\forall i\ \unsure{\mathrm{ob}\,I|\in\mathrm{Ob}\,I}$,
\scanpage{65}
\mspage{52}
% archive sheet number: 64

\inserted{considérant} le morphisme induit par $u$
\[
 I_{/i}\longrightarrow I'_{/i'}\qquad(\text{où }i'=u(i)),
\]
celui-ci est (\uline{globalement}) \unsure{ft|fort} cofiltrant, et que pour tt ouvert $U$ de $I_{/i}$, désignant par $U'$ l'ouvert engendré par son image, le morphisme induit $\unsure{v}:U\unsure{\mathbin{-}|\to}U'$ est encore \unsure{ft|fort} cofiltrant et a fortiori \uline{cofinal}.

\deleted{Considérons donc le diagramme commutatif, pour un $F'$ de $M(I'^{\unsure{\circ|0}})$ i.e. $F':I'\mathbin{-}M$}

\deleted{\diagram{\begin{tikzcd}[ampersand replacement=\&,column sep=large]
 \varinjlim_U u^*(F')\deleted{\unsure{(i)}}\arrow[r,"c"]\arrow[d,"{\alpha_{U,i}}"'] \& \varinjlim_{\unsure{V|U'}} F'\arrow[d,"{\alpha_{U',i'}}"]\\
 u^*(F')(i)\arrow[r,"{\sim\,\mathrm{Id}}"] \& F'(i')\ .
\end{tikzcd}}}

\deleted{Par cofinalité de $v$, la flèche can $c$ est iso, donc si \deleted{$\unsure{\alpha_{U',i'}}$} $\alpha_{U',i'}$ est iso, il en est de $=$ de $\alpha_{U,i}$. Donc si $F'$ est dans \deleted{$M(I'^\circ)_C$ \unsure{i.e.}}}

Soit donc $f':X'\mathbin{-}Y'$ une flèche \deleted{\unsure{dans}} $\unsure{\in} C_{I'^\circ}$, prouvons que \deleted{$\unsure{u^*(f')}$}
\[
 \underbrace{u^*(f')}_{f}:\underbrace{u^{\unsure{*|'}}(X')}_{X}\mathbin{-}\underbrace{u^{\unsure{'|*}}(Y')}_{Y}
\]
est dans $C_{I^\circ}$, i.e. que pour tt $i\in\deleted{\unsure{\mathrm{Ob}}} I$, et $U$ ouvert de $I_{/i}$, dans le diagramme canonique,

\diagram{\begin{tikzcd}[ampersand replacement=\&,column sep=large,row sep=small]
 \& X_U\arrow[r,dash]\arrow[d]\arrow[dr,phantom,"{\text{cocart.}}" description] \& X(i)\arrow[d]\\
 (\alpha) \& Y_U\arrow[r,dash] \& Y_{U,i}\arrow[d,"{v_{U,i}}"]\\
 \& \& Y\unsure{(i)}
\end{tikzcd}}
\scanpage{66}
\mspage{53}
% archive sheet number: 65

la flèche $v_{U,i}$ est $\unsure{\in|\text{ds}} C$. Le diagramme $(\alpha)$ s'envoie dans le diagramme similaire pour $f':X'\mathbin{-}Y'$, $i'=u(i)$, et l'ouvert $U'\subset I'_{/i'}$ engendré par $u(U)$ :

\diagram{\begin{tikzcd}[ampersand replacement=\&,column sep=small,row sep=small]
 \& X_U\arrow[rr,dash]\arrow[dr,"{f_U}"']\arrow[ddd,"{c^{X'}_{\deleted{\unsure{X}}U}}"'] \& \& X(i)\arrow[dr,dash]\arrow[ddd,"{\mathrm{id}}" near end]\arrow[dl,phantom,"{\text{cocart}}" description] \& \& \\
 \& \& Y_U\arrow[rr]\arrow[ddd,"{c^{Y'}_{\deleted{\unsure{Y}}U}}"'] \& \& Y_{U,i}\arrow[dr,"{v_{U,i}}"]\arrow[ddd,"{c^{f'}_{U,i}}"'] \& \\
 (\alpha\alpha) \& \& \& \& \& Y(i)\arrow[ddd,"{\mathrm{id}}"]\\
 \& X'_{U'}\arrow[rr]\arrow[dr,"{f'_{U}}"'] \& \& X'(i')\arrow[dr,dash]\arrow[dl,phantom,"{\text{cocart}}" description] \& \& \\
 \& \& Y'_{U'}\arrow[rr] \& \& Y'_{U',i'}\arrow[dr,"{v_{U',i'}}"'] \& \\
 \& \& \& \& \& Y'(i')\ .
\end{tikzcd}}

Par cofinalité, \deleted{\unsure{les}} deux flèches verticales $c^{X'}_U$ et $c^{Y'}_U$ sont iso, \deleted{\unsure{donc}} plus de $X(i)\overset{\mathrm{id}}{\mathbin{-}} X'(i')$, donc \deleted{$\unsure{Y_{U,i}}$} la 4\textsuperscript{e} flèche verticale du cube, $c^{f'}_{U,i}$, est iso. Comme $v_{U',i'}$ est $\in C$, $v_{U,i}$ aussi, on \unsure{gagne|a gagné}.

\deleted{\unsure{Prouvons} Nous n'avons pas utilisé le th. 3.1, mais nous allons l'utiliser pour prouver b).}

\inserted{Prouvons b).} Soit donc \deleted{$f$ \unsure{une}}
\[
 f:X\longrightarrow Y
\]
une flèche dans $C_I$, et soit
\[
 f'=u_!(f):u_!(X)=X'\longrightarrow u_!(Y)=Y',
\]
prouvons que c'est dans $C_{I'^{\unsure{\circ|0}}}$, donc que
\scanpage{67}
\mspage{54}
% archive sheet number: 66

pour deux ouverts $U'\subset V'$ de $I'$, \deleted{dans le \unsure{dia}} (où on peut supposer que $V'=I'_{/i'}$), dans le diagramme obtenu de $(\alpha\alpha)$, $v'_{U',V'}$ soit dans $C$. Soit $U=u^{-1}(U')$, $V=u^{-1}(V')$, \deleted{de sorte que par transit\unsure{.}}, et considérons le diagramme similaire à $(\alpha\alpha)$, défini par les deux inclusions d'ouverts $U'\subset V'$, $U\subset V$. Cette fois, comme la formation de $u_!$ \unsure{commute} aux restrictions aux ouverts, \inserted{par transitivité des $\underline{\lim}$} les trois flèches verticales $c^X_{U'}$, $c^Y_{U'}$, $c^X_{V'}$ sont iso, donc aussi $c^f_{U',V'}$.

De $=$ la flèche $Y_V\mathbin{-}Y'_{V'}$ (à la place de $\unsure{Y|X}(i)\mathbin{-}Y'(i')$) est iso, Comme $v_{U,V}\in C$, il en est de $=$ de $v_{U',V'}$, ce qui prouve que $f'\in C_{I'^{\unsure{\circ|0}}}$. \deleted{Il \unsure{en sera}} On \unsure{procède} de $=$ pour $TC_I$. Enfin, l'assertion que $u_!$ applique $M(I^\circ)_C$ dans $M(I'^\circ)_{C\unsure{*}}$ résulte des \deleted{fait} définitions et de $u_!(C_{I^\circ})\subset C_{I'^\circ}$.

\uline{Corollaire} 3.3. (de Prop. 3.2. et Th. 3.1)

Soit $u:I\longrightarrow I'$ comme dans 3.2. Supposons de plus $I$ \deleted{\unsure{locm}} artinien, et
\scanpage{68}
\mspage{55}
% archive sheet number: 67

qu'on soit \inserted{pour $M,W,C,I$} dans l'une des conditions a), b) du th. 3.1. Alors $u_!$ induit une flèche

\[
 u_{!\,\mathrm{cof}}:M(I^{\unsure{\circ|0}})_C\mathbin{-} M(I'^\circ)_C\subset M(I'^\circ),
\]

qui transforme flèches de $W(I^\circ)$ en flèches de $W(I'^\circ)$ (i.e compatible avec les localiseurs), Il induit donc un foncteur sur les localisées, d'où \deleted{par} \inserted{un} \deleted{$I'$ satisfaisant \unsure{égalt}} foncteur \deleted{foncteur}

\[
 u_!^W:\mathrm{Ho}_W(I^\circ)\longrightarrow\mathrm{Ho}_W(I'^\circ).\tag*{(3.3.1)}
\]

% slanted
\margin{NB Cette flèche est définie, $=$ sans supposer $I'$ ordonné, car $u_{!\,\mathrm{cof}}$ transforme \unsure{une flèche|flèches} de $W(I^\circ)$ en $W(I'^\circ)$. Cf. \unsure{aussi|mém.|même} p.~\unsure{86|88|56}}

\uline{Dém}. Par le th. 3.1, on sait que $(M(I^\circ),W(I^\circ),C\deleted{\unsure{(I^\circ)}}_{I^\circ})$ est une structure de cat. de modèles à cofibration. On sait par ailleurs que $TC_{I^\circ}=C_{I^\circ}\cap W(I^\circ)$, par Th. 2.1. $3^\circ$ (p.~27). Donc par 3.\unsure{2|3}, $u_!$ applique $C_{I^\circ}\cap W(I^\circ)$ dans $TC_{I'^\circ}\subset C_{I'^\circ}\cap\unsure{W(I'^\circ)}$, et a fortiori, en des flèches $\in W(I'^\circ)$.

Par ailleurs, par la théorie des cat à cofibrations, on trouve que toute flèche \inserted{$f$} \deleted{de $W(I^\circ)$ qui est dans} entre objets cofibrants est composée $f=pi$, où $i\in C_{I^\circ}$, et \deleted{deux flèches qui \unsure{est}} où $p$ est $\in W(I^\circ)$ \inserted{et de plus} \uline{inverse à gauche} d'une flèche de $C_{I^\circ}\cap W(I^\circ)$, \unsure{\deleted{de la}|donc} d'une flèche de $(TC)_{I^\circ}$. Alors
